\scanpage{243}
\[
\operatorname{Pic}(X_n)\overset{\sim}{\longrightarrow}\underline{\operatorname{Pic}}_{X_n/k}(k)
\]
et du th. précédent.

Dans le cas général, Notons \deleted{\unsure{d'abord}} ceci [\deleted{\unsure{si|où}} $k'/k$ est une extension
finie de $k$ \unsure{(G)}, \unsure{qu'on|bien} \unsure{choisira|choisie} de façon

\deleted{\uline{Lemme} \unsure{Sous|Sans} les conditions du corollaire
[car. 0 \unsure{essentielle|considérée}, mais on \unsure{le rappelle|la rappelle}],
le \unsure{conoyau|noyau} de $\operatorname{Pic}(X_{m})\unsure{-|\longrightarrow}\operatorname{Pic}(X_n)$
(pour $m\geq n$) est sans torsion.

On est ramené au cas où $m=n+1$,
où cela \deleted{\unsure{se}} résulte du fait que ce
conoyau est $\subset H^2(X_0,\deleted{\unsure{\mathscr F_n|F_n}}\inserted{F^{(n)}})$, $F^{(n)}_{\deleted{\unsure{n}}}$
un Module \unsure{coh|coh.} convenable sur $X_0$.}

que \deleted{\unsure{les|le}} $X_0\otimes_k k'$ \deleted{\unsure{soit}}satisfasse à la condition
énoncée plus haut] : on a une suite exacte

% slanted
\margin{\uline{\unsure{référence|références|vérifions}}?? \uline{À expliciter}}
\[
0\longrightarrow\operatorname{Pic}(X_n)\longrightarrow\underline{\operatorname{Pic}}_{X_n/k}(k)\longrightarrow\check H^2(k'/k,M_n)
\]
où $M_n$ est le schéma en groupes [comm. lisses
connexes] \deleted{\unsure{tel|tels}} qui décrit les unités
de la $k$-algèbre \deleted{\unsure{$H$}} $H^0(X_n,\mathcal O_{X_n})$,

On note que le système projectif des
$(M_n)$ satisfait M.L. [groupes \deleted{\unsure{alg. connexes}}\inserted{\uline{de type fini} sur $k$ !}]
\unsure{et est|le est} : noyaux \deleted{\unsure{unipotents}} \unsure{finis|lisses|tous}, \uline{vectoriels}\ldots]
% archive sheet number: 242
